\section{The branch inequality: topology and weights}
\label{sec:branch-topology}

The negative term in \eqref{eq:difficulty-local-deficit} counts excess
normalization branches along surfaces of the varying boundary.  We now
bound this branch excess by divisor-cycle ranks on the normalizations
and valuations centered on singular curves.  The estimate applies to
any terminal pair and birational morphism satisfying the hypotheses
below.

The proof has two parts.  Normalization relations and the comparison
with intersection cohomology first bound the excess by divisor-cycle
ranks and punctured degree-two cohomology along singular curves.  In
Section~\ref{sec:local-valuations}, classes in that local cohomology
fixed by monodromy produce distinct divisorial valuations of the original
function field.  Those valuations supply the remaining term of the
estimate.

All cohomology and homology in these two sections have rational
coefficients and refer to the complex analytic topology.  We use the
perverse normalization of intersection complexes: if $Y$ is smooth of
dimension four, then $\mathrm{IC}_Y=\mathbb Q_Y[4]$.

\begin{theorem}[Branch inequality]\label{thm:branch}
Let $p\colon Y\to X$ be a projective birational morphism of normal
projective fourfolds over $\mathbb C$.  Suppose that $Y$ is
$\mathbb Q$-factorial and terminal and that $(Y,\Theta)$ is an effective
terminal pair with rational boundary.  Let
\[
 Q=\bigcup_{j=1}^{m}Q_j
\]
be a reduced union of distinct positive-coefficient components of
$\Theta$, and assume $\dim p(Q)\leq 1$.
Write $Q_j^\nu$ for the normalization of $Q_j$.  For each irreducible
surface $V\subset Q$, let $r_V$ be the number of prime divisors of
$\coprod_jQ_j^\nu$ mapping onto $V$.  Then
\begin{equation}\label{eq:branch-inequality}
 \begin{aligned}
 \sum_{V\subset Q}(r_V-1)
 &\leq\sum_{j=1}^{m}\rho_*(Q_j^\nu)\\
 &\quad+\#\bigl\{v:\operatorname{codim}_Yc_Y(v)=3,\;
          c_Y(v)\subset Q,\;a(v;Y,\Theta)<2\bigr\}.
 \end{aligned}
\end{equation}
Here the valuations are normalized divisorial orders, and the sum on
the left has only finitely many nonzero terms.
\end{theorem}

The branches in this statement are algebraic prime divisors.
In particular, an irreducible divisor mapping to $V$ with degree
greater than one counts as one branch; its degree will occur in the
pushforward map.  No crepancy assumption on $p$ is made.

Set
\begin{equation}\label{eq:branch-loci}
 C=p(Q)_{\mathrm{red}},\qquad
 P=(p^{-1}C)_{\mathrm{red}}.
\end{equation}
Let $P_1$ be the reduced union of the three-dimensional components
of $P$.  Thus $Q\subset P_1\subset P$, and $P\setminus P_1$ has complex
dimension at most two.  If $Q$ is empty, the theorem is immediate; in
what follows we may assume otherwise.

\subsection{Normalization relations}

For a projective variety $R$, its homology carries the mixed Hodge
structure dual to its cohomology.  In particular the top homology
$H_4(F)$ of a projective variety of dimension at most two is freely
generated by the fundamental classes of its irreducible surfaces and
is pure of weight $-4$.  We write $W$ for the weight filtration.

\begin{lemma}[Conductor estimate]\label{lem:conductor}
With the notation above,
\begin{equation}\label{eq:conductor-rank}
 \sum_{V\subset Q}(r_V-1)-\sum_j\rho_*(Q_j^\nu)
 \leq
 \operatorname{rank}\operatorname{Gr}^{W}_{4}
       \bigl[H^5(P)\longrightarrow H^5(Q)\bigr].
\end{equation}
\end{lemma}

\begin{proof}
For a reduced projective pure threefold $R$, let
$\nu_R\colon N_R\to R$ be its normalization.  Choose a closed reduced
subset $A_R\subset R$ of dimension at most two such that $\nu_R$ is an
isomorphism away from $A_R$, and put
$F_R=(\nu_R^{-1}A_R)_{\mathrm{red}}$.  The normalization square is a
proper excision square.  Its homology sequence, compatible with mixed
Hodge structures, contains
\begin{equation}\label{eq:normalization-excision}
 H_5(R)\xrightarrow{\partial_R} H_4(F_R)
 \longrightarrow H_4(A_R)\oplus H_4(N_R).
\end{equation}
This follows by comparing the localization sequences for
$(R,A_R)$ and $(N_R,F_R)$, whose open complements are isomorphic;
the mixed Hodge structures and their functoriality are those of
\cite[Proposition~8.2.2 and Proposition~8.3.9]{Deligne1974}.

Apply this construction first to $R=P_1$, taking $A_{P_1}$ to contain
the nonisomorphism locus of its normalization.  For $R=Q$ take
$A_Q=Q\cap A_{P_1}$.  The normalization $N_Q=\coprod_j Q_j^\nu$ is a
union of entire components of $N_{P_1}$, so these choices give a
morphism between the two normalization squares.

Let $V$ run over the surface components of $A_Q$.  Write
$E_{V,1},\ldots,E_{V,r_V}$ for the surface components of $F_Q$ over
$V$, and let $e_{V,\alpha}=[\mathbb C(E_{V,\alpha}):\mathbb C(V)]$.
On their fundamental classes, the first component of the map in
\eqref{eq:normalization-excision} is
\[
 \bigoplus_{\alpha=1}^{r_V}\mathbb Q[E_{V,\alpha}]
 \longrightarrow \mathbb Q[V],\qquad
 \sum_\alpha t_\alpha[E_{V,\alpha}]
 \longmapsto \Bigl(\sum_\alpha e_{V,\alpha}t_\alpha\Bigr)[V].
\]
Its kernel has dimension $r_V-1$.  Every surface with $r_V>1$ lies
in $A_Q$, because normalization is an isomorphism elsewhere.
Consequently the kernel $B$ of $H_4(F_Q)\to H_4(A_Q)$ satisfies
\[
 \dim B=\sum_{V\subset Q}(r_V-1).
\]
All these terms are finite in number.  The image of
$B\to H_4(N_Q)$ lies in the span of divisor classes, whose dimension
is $\sum_j\rho_*(Q_j^\nu)$.  Therefore
\[
 U:=\ker\bigl[B\longrightarrow H_4(N_Q)\bigr],\qquad
 \dim U\geq\sum_{V\subset Q}(r_V-1)-\sum_j\rho_*(Q_j^\nu).
\]
By exactness, $U$ is the image of $\partial_Q$.

The map $H_4(F_Q)\to H_4(F_{P_1})$ is injective: each surface
generator of the source is the same surface in one of the selected
normalization components of $P_1$, and different such generators
remain distinct.  Naturality of the boundary maps now shows that
$\partial_{P_1}$ maps the image of $H_5(Q)\to H_5(P_1)$ onto a
space containing a copy of $U$.  Since $H_4(F_{P_1})$ is pure of
weight $-4$, strictness of morphisms of mixed Hodge structures
\cite[Th\'eor\`eme~2.3.5(iii)]{Deligne1971} gives
\[
 \dim U\leq
 \operatorname{rank}\operatorname{Gr}^{W}_{-4}
          \bigl[H_5(Q)\longrightarrow H_5(P_1)\bigr].
\]
Finally the localization sequence in Borel--Moore homology contains
\[
 H^{\mathrm{BM}}_6(P\setminus P_1)
 \longrightarrow H_5(P_1)\longrightarrow H_5(P).
\]
Its first term vanishes by the dimension bound, so its second arrow
is injective.  Duality between homology and cohomology, with weight
sign reversed, proves \eqref{eq:conductor-rank}.
\end{proof}

\subsection{A purity statement over a curve}

The conductor estimate has reduced the problem to the weight-four
part of $H^5(P)\to H^5(Q)$.  We next prove that
$\mathbb H^1(P,\mathrm{IC}_Y|_P)$ is pure of weight five.  Comparison
with ordinary cohomology will then isolate the weight-four contribution
on singular curves.  The decomposition theorem is applied before
restriction to $C$;
the support estimates below identify which of its summands can
contribute in degree one.

\begin{lemma}[Intersection-complex restriction]\label{lem:restriction-purity}
Let $p\colon Y\to X$ be a projective birational morphism of normal
projective fourfolds over $\mathbb C$, and let $C\subset X$ be a
closed reduced subset of dimension at most one.  Put
$P=(p^{-1}C)_{\mathrm{red}}$.  Then
\[
 \mathbb H^1(P,\mathrm{IC}_Y|_P)
\]
is pure of weight five.
\end{lemma}

\begin{proof}
We use intersection complexes and direct images in the derived
category of mixed Hodge modules, and use the same symbols for their
underlying rational complexes.  The intersection complex of the
irreducible fourfold $Y$ is pure of weight four.  The projective
decomposition theorem gives a finite direct sum
\begin{equation}\label{eq:ic-decomposition}
 Rp_*\mathrm{IC}_Y\simeq
 \mathrm{IC}_X\oplus\bigoplus_{(Z,i)}M_{Z,i}[-i],
\end{equation}
where each $Z$ is a proper irreducible closed subset of $X$, and
$M_{Z,i}$ is a pure Hodge module of strict support $Z$ and weight
$4+i$.  Multiple summands with the same support and index are
allowed.  Birationality gives the unique full-support summand
$\mathrm{IC}_X$ in degree zero.  The weight and decomposition
statements used here are \cite[\S\S1.6--1.8, 1.12--1.13]{Saito1989}.

We claim that every proper-support summand in
\eqref{eq:ic-decomposition} satisfies
\begin{equation}\label{eq:proper-support-bound}
 i+z\leq 2,\qquad z=\dim Z.
\end{equation}
Choose an algebraic stratification adapted to $\mathrm{IC}_Y$ and
refine it by $p^{-1}Z$.  Also stratify the target so that over a
dense smooth open subset $Z^\circ$ of $Z$ every source stratum
meeting a fibre has constant fibre dimension; strata whose images
are not dense in $Z$ are avoided by shrinking $Z^\circ$.
This is possible by constructibility and generic stratified
triviality.  In particular it includes the strata introduced by
fibre-dimension jumps before shrinking.

Let $B\subset p^{-1}Z$ be one of the refined strata dominating
$Z$, and put $u=\dim B$.  Since $p$ is birational and $Z$ is
proper, $p^{-1}Z$ is a proper closed subset of $Y$, and hence
$u\leq3$.  On $B$ we have
\begin{equation}\label{eq:source-stalk-bound}
 \mathcal H^q(\mathrm{IC}_Y)|_B=0\quad\text{for }q>-u-1.
\end{equation}
Indeed, if $B$ lies in the smooth open stratum, the only degree is
$-4\leq-u-1$.  Otherwise let $t\geq u$ be the dimension of the
original stratum containing it.  The strict support condition for
an intersection complex gives vanishing for $q>-t-1$, and this
implies \eqref{eq:source-stalk-bound}.  This is the stalk support
condition in the middle-extension axioms
\cite[\S2.3 and \S3.3, axiom {\rm[AX1]}(c)]{GM1983}.

For $x\in Z^\circ$, the stratum $B\cap p^{-1}(x)$ has dimension
$u-z$.  Compactly supported cohomology on this stratum vanishes
above degree $2(u-z)$.  By \eqref{eq:source-stalk-bound} its
contribution to the cohomology of the fibre with coefficients in
$\mathrm{IC}_Y$ vanishes in degrees greater than
\[
 -u-1+2(u-z)=u-2z-1\leq2-2z.
\]
A finite filtration by strata and proper base change give the
same upper bound for the stalk of $Rp_*\mathrm{IC}_Y$ at $x$.
On a sufficiently small $Z^\circ$, the nonzero stalk of
$M_{Z,i}[-i]$ is in degree $i-z$.  Since it is a direct summand,
$i-z\leq2-2z$, proving \eqref{eq:proper-support-bound}.

Consider now a support $Z$ not contained in $C$.  Stratify $Z\cap C$
compatibly with $M_{Z,i}$, and let $S$ be a stratum of dimension
$s$.  Then $s\leq z-1$.  The stalks of $M_{Z,i}$ on $S$ vanish
above degree $-s-1$: on the local-system open of $Z$ their only
degree is $-z\leq-s-1$, and elsewhere this follows from the
strict support condition, including after refinement as above.
Thus the largest possible degree of a contribution from $S$ to
the compactly supported cohomology of $M_{Z,i}[-i]|_C$ is
\[
 i-s-1+2s=i+s-1\leq i+z-2\leq0.
\]
Filtration by strata shows that this summand contributes nothing
to degree one on $C$.  For the full-support summand
$\mathrm{IC}_X$, each stratum of $C$ has dimension $s\leq1$ and
the same stalk argument gives the upper bound
$-s-1+2s=s-1\leq0$.  It too contributes nothing in degree one.

It remains to consider $Z\subset C$.  Such a support is projective,
and its contribution is
\[
 \mathbb H^{1-i}(Z,M_{Z,i}),
\]
which is pure of weight $(4+i)+(1-i)=5$ by projective direct-image
purity.  Proper base change identifies the degree-one cohomology
of \eqref{eq:ic-decomposition} restricted to $C$ with
$\mathbb H^1(P,\mathrm{IC}_Y|_P)$.  All its nonzero summands are
therefore pure of weight five, as claimed.
\end{proof}

\subsection{The comparison cone and singular curves}

Return to the hypotheses of Theorem~\ref{thm:branch}, and write
$j\colon Y_{\mathrm{reg}}\hookrightarrow Y$ for the smooth-locus
inclusion.  Terminality implies
$\dim\operatorname{Sing}Y\leq1$.  For every irreducible curve
$T\subset\operatorname{Sing}Y$, choose a connected smooth dense
open subset $T^\circ$ on which the complexes below have locally
constant cohomology.  We may remove further finitely many points
whenever needed.  Define
\begin{equation}\label{eq:curve-local-system}
 \mathcal V_T=(R^2j_*\mathbb Q_{Y_{\mathrm{reg}}})|_{T^\circ},
 \qquad
 b_T=\dim H_c^2(T^\circ,\mathcal V_T).
\end{equation}
Equivalently, $b_T$ is the dimension of the coinvariant space of
a fibre of $\mathcal V_T$ under its monodromy representation.
Removing finitely many further points does not change this number.

\begin{lemma}[Local-system bound]\label{lem:local-system-bound}
For $P$ and $Q$ in \eqref{eq:branch-loci},
\begin{equation}\label{eq:local-system-rank}
 \operatorname{rank}\operatorname{Gr}^{W}_{4}
       \bigl[H^5(P)\longrightarrow H^5(Q)\bigr]
 \leq
 \sum_{\substack{T\text{ irreducible curve}\\
                  T\subset\operatorname{Sing}Y\cap Q}}b_T.
\end{equation}
\end{lemma}

\begin{proof}
There is a comparison morphism, equal to the identity on
$Y_{\mathrm{reg}}$, in the derived category of mixed Hodge modules
\cite[\S1.14]{Saito1989}.  Let $\mathcal T$ be its cone:
\begin{equation}\label{eq:comparison-triangle}
 \mathbb Q_Y[4]\longrightarrow\mathrm{IC}_Y
 \longrightarrow\mathcal T\longrightarrow\mathbb Q_Y[5].
\end{equation}
The same comparison complex is studied by Park--Popa
\cite[Definition~6.1]{ParkPopaHodge} as their rational-homology-manifold
defect complex shifted by one.
Restricting to $P$ yields the exact sequence of mixed Hodge
structures
\begin{equation}\label{eq:comparison-cohomology}
 \mathbb H^0(P,\mathcal T|_P)
 \xrightarrow{\delta_P}H^5(P)
 \longrightarrow\mathbb H^1(P,\mathrm{IC}_Y|_P).
\end{equation}
Lemma~\ref{lem:restriction-purity} makes the last term pure of
weight five.  Consequently
\[
 W_4H^5(P)\subset\operatorname{im}\delta_P.
\]
Naturality under $Q\subset P$ implies that the image of
$W_4H^5(P)$ in $H^5(Q)$ is contained in the image of
\[
 \delta_Q\colon\mathbb H^0(Q,\mathcal T|_Q)\longrightarrow H^5(Q).
\]
Strictness of the restriction map for the weight filtration
therefore gives
\begin{equation}\label{eq:cone-rank}
 \operatorname{rank}\operatorname{Gr}^{W}_{4}
       \bigl[H^5(P)\longrightarrow H^5(Q)\bigr]
 \leq\dim\mathbb H^0(Q,\mathcal T|_Q).
\end{equation}

The cone is supported on $\operatorname{Sing}Y$.  On a stratum
of this locus of dimension $s\in\{0,1\}$, the strict stalk bound
for $\mathrm{IC}_Y$ and the triangle
\eqref{eq:comparison-triangle} give
\[
 \mathcal H^q(\mathcal T)=0\qquad(q>-s-1).
\]
In particular a point stratum has no contribution in degree zero.
On $T^\circ$ the only possible contribution in total degree zero
is
\[
 H_c^2(T^\circ,\mathcal H^{-2}(\mathcal T)|_{T^\circ}).
\]
At the generic codimension-three stratum, the middle-extension
construction retains the cohomology of $Rj_*\mathbb Q$ through
unshifted degree two.  Thus, by the attaching truncation in
\cite[\S3.1 and \S3.3, axiom {\rm[AX1]}(d)]{GM1983},
\begin{equation}\label{eq:cone-stalk}
 \mathcal H^{-2}(\mathcal T)|_{T^\circ}
 =\mathcal H^{-2}(\mathrm{IC}_Y)|_{T^\circ}
 =\mathcal V_T.
\end{equation}
The first equality also follows directly from
\eqref{eq:comparison-triangle}, since the constant complex has
stalk cohomology only in degree $-4$.

Take the $T^\circ$ to be disjoint and to omit every intersection
point of singular curves.  The intersection of a singular curve
not contained in $Q$ with $Q$ consists of finitely many points.
Hence the complement in $\operatorname{Sing}Y\cap Q$ of the
$T^\circ$ for curves $T\subset Q$ is finite.  The open--closed
cohomology sequence and the point-stratum vanishing give a
surjection
\[
 \bigoplus_{T\subset\operatorname{Sing}Y\cap Q}
 \mathbb H_c^0(T^\circ,\mathcal T|_{T^\circ})
 \longrightarrow\mathbb H^0(Q,\mathcal T|_Q).
\]
The compact-support spectral sequence on a curve, together with
the stalk bound and \eqref{eq:cone-stalk}, identifies its summands
with $H_c^2(T^\circ,\mathcal V_T)$.  The dimension bound from this
surjection and \eqref{eq:cone-rank} prove the lemma.
\end{proof}

\begin{proof}[Proof of Theorem~\ref{thm:branch}]
Lemmas~\ref{lem:conductor} and \ref{lem:local-system-bound} give
\begin{equation}\label{eq:branch-before-valuations}
 \sum_{V\subset Q}(r_V-1)
 \leq\sum_j\rho_*(Q_j^\nu)
   +\sum_{T\subset\operatorname{Sing}Y\cap Q}b_T.
\end{equation}
Proposition~\ref{prop:local-valuations}, proved in the next section,
constructs for each such curve $T$ at least $b_T$ distinct
normalized divisorial valuations $v$ satisfying
\[
 c_Y(v)=T,\qquad a(v;Y,0)=2,\qquad a(v;Y,\Theta)<2.
\]
Valuations arising from different curves have different centers
and are distinct.  Lemma~\ref{lem:echoes} ensures that the set of
valuations of discrepancy below two with codimension-three center
is finite.  Substituting their number for the last sum in
\eqref{eq:branch-before-valuations} gives
\eqref{eq:branch-inequality}.
\end{proof}
